% Lean module: MonadFormal.Information.Testing.Hellinger.Variation
\paragraph{Comparison with operational total variation.}

\begin{theorem}[Density formula and Hellinger bounds]
\label{thm:hellinger-tv}
Let $\delta=d_{\mathrm{TV}}(P_r,P_s)$ use the normalized
measurable-test distance of \Cref{sec:simulation}. Then
\[
 \delta=\frac12\int|r^2-s^2|\,d\mu,
 \qquad
 H^2(P_r,P_s)\le\delta,
 \qquad
 \delta^2\le H^2(P_r,P_s)\bigl(2-H^2(P_r,P_s)\bigr).
\]
The density formula is an equality with that operational distance,
not an independent definition.
\end{theorem}
\begin{proof}
Put $f=r^2-s^2$. It is integrable and $\int f\,d\mu=0$.
For a measurable event $A$, its complement has the opposite signed
integral. Applying the integral triangle inequality on both pieces gives
$2|\int_A f\,d\mu|\le\int|f|\,d\mu$.
Equality is attained by $A=\{f\ge0\}$, since the positive and negative
parts have equal integral. The event characterization of total variation
therefore proves the density formula.

For nonnegative $r,s$, $(r-s)^2\le|r^2-s^2|$ pointwise, giving
the lower bound. For the upper bound, factor the density difference
and apply Cauchy--Schwarz:
\[
 \delta=\frac12\int|(r-s)(r+s)|\,d\mu
 \le\frac12\|r-s\|_2\|r+s\|_2.
\]
Now $\|r-s\|_2^2=2H^2$ and
$\|r+s\|_2^2=2+2\rho=4-2H^2$.
Squaring the nonnegative bound proves the last inequality.
\leanproof{MonadFormal.Information.abs_setIntegral_le_half_integral_abs,MonadFormal.Information.integral_abs_mul_le_L2_norm,MonadFormal.Information.RootDensity.density_diff_integrable,MonadFormal.Information.RootDensity.density_diff_integral,MonadFormal.Information.RootDensity.tv_eq_density_integral,MonadFormal.Information.RootDensity.tv_toReal_eq_density_integral,MonadFormal.Information.RootDensity.hellingerSq_le_tv,MonadFormal.Information.RootDensity.tv_sq_le_hellingerSq}
\end{proof}
